\documentclass[10pt,a4paper]{article}

\usepackage[utf8]{inputenc}
\usepackage[T1]{fontenc}
\usepackage[french]{babel}
\usepackage{amsmath,amssymb,amsthm}
\usepackage[top=1cm,bottom=1.2cm,left=1.5cm,right=1.5cm]{geometry}
\usepackage{hyperref}
\hypersetup{colorlinks=true,linkcolor=black,citecolor=blue,urlcolor=blue}

\newtheorem{proposition}{Proposition}

\title{Rétablissement de l'homogénéité : la statistique culminante}
\author{Salomé Fonvielle}
\date{}

\begin{document}
\maketitle

\noindent\emph{Question traitée.} Montrer que la statistique culminante du CUSUM
s'écrit $S_{\max}(H) = \max_{0 \le k \le j \le H}\bigl[A(k,j) - (j-k)\,\delta_P\bigr]$,
le maximum portant sur toutes les sous-fenêtres $]k,j] \subseteq [0,H]$.

\paragraph{Notations.} Les pas post-rupture sont indexés $t = 1,\dots,H$, l'instant
$t=0$ étant la rupture elle-même. On pose, conformément à l'énoncé,
\[
  A(k,j) \;=\; \sum_{t=k+1}^{j} \bigl(e_t - \hat p_0\bigr),
  \qquad
  X_t \;=\; e_t - \hat p_0 - \delta_P,
  \qquad
  S_0 = 0,
  \qquad
  S_t = \max\bigl(0,\, S_{t-1} + X_t\bigr),
\]
et $S_{\max}(H) = \max_{0 \le t \le H} S_t$. Le socle est ici le socle \emph{estimé}
$\hat p_0$, comme dans toute la question~C. Par convention une somme vide vaut $0$,
de sorte que $A(j,j) = 0$.

\begin{proposition}[Forme close de Lindley]\label{prop:lindley}
Pour tout $j \in \{0,\dots,H\}$,
\begin{equation}
  S_j \;=\; \max_{0 \le k \le j}\ \sum_{t=k+1}^{j} X_t
  \;=\; \max_{0 \le k \le j}\ \bigl[A(k,j) - (j-k)\,\delta_P\bigr].
  \label{eq:lindley}
\end{equation}
En passant au maximum sur $j$, il vient
$S_{\max}(H) = \max_{0 \le k \le j \le H}\bigl[A(k,j) - (j-k)\,\delta_P\bigr]$.
\end{proposition}

\begin{proof}
Par récurrence sur $j$.

\smallskip
\noindent\emph{Initialisation.} Pour $j = 0$, le seul indice admissible est $k = 0$ et
la somme correspondante est vide, donc nulle : le membre de droite vaut $0 = S_0$.

\smallskip
\noindent\emph{Hérédité.} Supposons~\eqref{eq:lindley} vraie au rang $j-1$. Alors
\[
  S_j \;=\; \max\bigl(0,\; S_{j-1} + X_j\bigr)
      \;=\; \max\Bigl(0,\; \max_{0 \le k \le j-1} \sum_{t=k+1}^{j-1} X_t \;+\; X_j\Bigr)
      \;=\; \max\Bigl(0,\; \max_{0 \le k \le j-1} \sum_{t=k+1}^{j} X_t\Bigr),
\]
la dernière égalité résultant de ce que $X_j$ ne dépend pas de $k$ et peut donc entrer
dans le maximum. Or le terme $0$ est exactement la somme vide obtenue pour $k = j$ :
les deux arguments du maximum extérieur se réunissent en un seul maximum sur
$0 \le k \le j$, ce qui donne~\eqref{eq:lindley} au rang $j$.

\smallskip
\noindent\emph{Conclusion.} La seconde écriture s'obtient en développant la somme :
\[
  \sum_{t=k+1}^{j} X_t
  \;=\; \sum_{t=k+1}^{j} \bigl(e_t - \hat p_0\bigr) \;-\; (j-k)\,\delta_P
  \;=\; A(k,j) - (j-k)\,\delta_P ,
\]
puisque $\delta_P$ est retranché une fois par pas et que la sous-fenêtre $]k,j]$ en
compte $j-k$. Enfin
$S_{\max}(H) = \max_{0 \le j \le H} S_j
 = \max_{0 \le j \le H} \max_{0 \le k \le j} [\,\cdot\,]
 = \max_{0 \le k \le j \le H} [\,\cdot\,]$.
\end{proof}

\paragraph{Ce que l'identité rend visible.} La remise à zéro n'est pas une correction
technique : elle est \emph{équivalente} à la liberté de choisir le point de départ. Un
CUSUM qui repart de zéro à chaque fois que le compteur passerait sous la barre ne fait
rien d'autre que retenir, à chaque instant, la meilleure sous-fenêtre récente. C'est
pourquoi le maximum porte sur les couples $(k,j)$ et non sur les seuls $j$ : le
détecteur ne mesure pas la preuve accumulée depuis la rupture, mais la preuve accumulée
sur la sous-fenêtre la plus favorable.

Deux conséquences immédiates. D'abord $S_{\max}(H) \ge 0$, puisque le choix $k=j$ est
toujours disponible. Ensuite, et surtout, l'identité est \emph{déterministe} : elle vaut
sur chaque exécution prise isolément, sans moyenne ni hypothèse de loi. C'est ce qui
autorise le certificat de la question suivante.

\paragraph{Contrôle numérique.} La forme close~\eqref{eq:lindley} a été confrontée à la
récurrence naïve sur l'ensemble des trajectoires de la campagne : écart maximal
$1{,}07 \times 10^{-10}$, soit la précision machine. Le script d'analyse refuse de
poursuivre si ce contrôle échoue.

\end{document}
